%%% Hlavní soubor. Zde se definují základní parametry a odkazuje se na ostatní
%%% části. %%%

%% Verze pro jednostranný tisk:
% Okraje: levý 40mm, pravý 25mm, horní a dolní 25mm
% (ale pozor, LaTeX si sám přidává 1in)
\documentclass[12pt,a4paper]{report}
\setlength\textwidth{145mm}
\setlength\textheight{247mm}
\setlength\oddsidemargin{15mm}
\setlength\evensidemargin{15mm}
\setlength\topmargin{0mm}
\setlength\headsep{0mm}
\setlength\headheight{0mm}
% \openright zařídí, aby následující text začínal na pravé straně knihy
\let\openright=\clearpage

%% Pokud tiskneme oboustranně:
% \documentclass[12pt,a4paper,twoside,openright]{report}
% \setlength\textwidth{145mm}
% \setlength\textheight{247mm}
% \setlength\oddsidemargin{15mm}
% \setlength\evensidemargin{0mm}
% \setlength\topmargin{0mm}
% \setlength\headsep{0mm}
% \setlength\headheight{0mm}
% \let\openright=\cleardoublepage

%% Pokud používáte csLaTeX (doporučeno):
%\usepackage{czech}
%% Pokud nikoliv:
\usepackage[czech,english]{babel}
\selectlanguage{english}
\usepackage[T1]{fontenc}

%% Použité kódování znaků: obvykle latin2, cp1250 nebo utf8:
\usepackage[utf8]{inputenc}

%% Ostatní balíčky
\usepackage{graphicx}
\usepackage{amsthm}

%% Balíček hyperref, kterým jdou vyrábět klikací odkazy v PDF,
%% ale hlavně ho používáme k uložení metadat do PDF (včetně obsahu).
%% POZOR, nezapomeňte vyplnit jméno práce a autora.
\usepackage[ps2pdf,unicode]{hyperref}   % Musí být za všemi ostatními balíčky
\hypersetup{pdftitle=Fast and Trainable Tokenizer for Natural Languages}
\hypersetup{pdfauthor=Jiří Maršík}

%%% Drobné úpravy stylu

% Tato makra přesvědčují mírně ošklivým trikem LaTeX, aby hlavičky kapitol
% sázel příčetněji a nevynechával nad nimi spoustu místa. Směle ignorujte.
\makeatletter
\def\@makechapterhead#1{
  {\parindent \z@ \raggedright \normalfont
   \Huge\bfseries \thechapter. #1
   \par\nobreak
   \vskip 20\p@
}}
\def\@makeschapterhead#1{
  {\parindent \z@ \raggedright \normalfont
   \Huge\bfseries #1
   \par\nobreak
   \vskip 20\p@
}}
\makeatother

% Toto makro definuje kapitolu, která není očíslovaná, ale je uvedena v obsahu.
\def\chapwithtoc#1{
\chapter*{#1}
\addcontentsline{toc}{chapter}{#1}
}

\usepackage{mystyle}

\begin{document}

% Trochu volnější nastavení dělení slov, než je default.
\lefthyphenmin=2
\righthyphenmin=2


%%% Titulní strana práce

\pagestyle{empty}
\begin{center}

\large

Charles University in Prague

\medskip

Faculty of Mathematics and Physics

\vfill

{\bf\Large BACHELOR THESIS}

\vfill

\centerline{\mbox{\includegraphics[width=60mm]{img/logo.eps}}}

\vfill
\vspace{5mm}

{\LARGE Jiří Maršík}

\vspace{15mm}

% Název práce přesně podle zadání
{\LARGE\bfseries Fast and Trainable Tokenizer \\ for Natural Languages}

\vfill

% Název katedry nebo ústavu, kde byla práce oficiálně zadána
% (dle Organizační struktury MFF UK)
Institute of Formal and Applied Linguistics

\vfill

\begin{tabular}{rl}

Supervisor of the bachelor thesis: & RNDr. Ondřej Bojar, Ph.D. \\
\noalign{\vspace{2mm}}
Study program: & Computer Science \\
\noalign{\vspace{2mm}}
Specialization: & General Computer Science \\
\end{tabular}

\vfill

% Zde doplňte rok
Prague 2011

\end{center}

\newpage

%%% Následuje vevázaný list -- kopie podepsaného "Zadání bakalářské práce".
%%% Toto zadání NENÍ součástí elektronické verze práce, nescanovat.

%%% Na tomto místě mohou být napsána případná poděkování (vedoucímu práce,
%%% konzultantovi, tomu, kdo zapůjčil software, literaturu apod.)

\openright

\noindent
Dedicated to the work of B\'ela Tarr.

\newpage

%%% Strana s čestným prohlášením k bakalářské práci

\vglue 0pt plus 1fill

\noindent
I declare that I carried out this bachelor thesis independently, and only with
the cited sources, literature and other professional sources.

\medskip\noindent
I understand that my work relates to the rights and obligations under the Act
No. 121/2000 Coll., the Copyright Act, as amended, in particular the fact that
the Charles University in Prague has the right to conclude a license agreement
on the use of this work as a school work pursuant to Section 60 paragraph 1 of
the Copyright Act.

\vspace{10mm}

\hbox{\hbox to 0.5\hsize{%
In ........ date ............
\hss}\hbox to 0.5\hsize{%
signature
\hss}}

\vspace{20mm}
\newpage

%%% Povinná informační strana bakalářské práce

\vbox to 0.5\vsize{
\setlength\parindent{0mm}
\setlength\parskip{3mm}

\begin{singlespace}
\begin{otherlanguage}{czech}

Název práce:
Rychlý a trénovatelný tokenizér pro přirozené jazyky
% přesně dle zadání

Autor:
Jiří Maršík

Katedra:  % Případně Ústav:
Ústav formální a aplikované lingvistiky
% dle Organizační struktury MFF UK

Vedoucí bakalářské práce:
RNDr. Ondřej Bojar Ph.D.
% dle Organizační struktury MFF UK, případně plný název pracoviště mimo MFF UK

Abstrakt:
% abstrakt v rozsahu 80-200 slov; nejedná se však o opis zadání bakalářské
% práce
V této práci představujeme systém pro dezambiguaci hranic mezi tokeny a větami.
Charakteristickým znakem programu je jeho značná konfigurovatelnost a
všestrannost, tokenizér si dokáže poradit např. i s nepřerušovaným čínským
textem. Tokenizér používá klasifikátory založené na modelech s maximální
entropií a jedná se tudíž o systém strojového učení, kterému je nutné předložit
již tokenizovaná data k trénování. Program je doplněn nástrojem pro hlášení
úspěšnosti tokenizace, což pomáhá zejména při rychlém vývoji a ladění
tokenizačního procesu. Systém byl vyvinut pouze za pomoci multiplatformních
knihoven a při vývoji byl kladen důraz zejména na efektivitu a správnost. Po
nezbytném přehledu jiných tokenizérů a krátkém úvodu do teorie modelů s
maximální entropií se většina textu práce zabývá vlastní implementací
tokenizéru a vyhodnocením jeho úspěšnosti.

Klíčová slova:
tokenizace, segmentace, maximální entropie, předzpracování textu
% 3 až 5 klíčových slov

\end{otherlanguage}
\end{singlespace}

\vss}\nobreak\vbox to 0.49\vsize{
\setlength\parindent{0mm}
\setlength\parskip{3mm}

\begin{singlespace}

Title:
Fast and Trainable Tokenizer for Natural Languages
% přesný překlad názvu práce v angličtině

Author:
Jiří Maršík

Department:
Institute of Formal and Applied Linguistics
% dle Organizační struktury MFF UK v angličtině

Supervisor:
RNDr. Ondřej Bojar Ph.D.
% dle Organizační struktury MFF UK, případně plný název pracoviště
% mimo MFF UK v angličtině

Abstract:
In this thesis, we present a data-driven system for disambiguating token and
sentence boundaries. The implemented system is highly configurable and
versatile to the point its tokenization abilities are able to satisfactorily
segment unbroken Chinese text. The tokenizer relies on maximum entropy
classifiers and requires tokenized and segmented text as training data. The
program is accompanied by a tool for reporting the performance of the
tokenization which helps to rapidly develop and tune the tokenization process. The
system was built with multi-platform libraries only and with emphasis on speed
and correctness. After a necessary survey of other tools for text tokenization
and segmentation and a short introduction to maximum entropy modelling a large
part of the thesis focuses on the particular implementation we developed and
its evaluation.
% abstrakt v rozsahu 80-200 slov v angličtině; nejedná se však o překlad zadání
% bakalářské práce

Keywords:
% 3 až 5 klíčových slov v angličtině
tokenization, segmentation, maximum entropy, text preprocessing

\end{singlespace}

\vss}

\newpage


%%% Strana s automaticky generovaným obsahem bakalářské práce. U matematických
%%% prací je přípustné, aby seznam tabulek a zkratek, existují-li, byl umístěn
%%% na začátku práce, místo na jejím konci.

\openright
\pagestyle{plain}
\setcounter{page}{1}
\tableofcontents

%%% Jednotlivé kapitoly práce jsou pro přehlednost uloženy v samostatných
%%% souborech
\chapter*{Introduction}
\addcontentsline{toc}{chapter}{Introduction}


\chapter{Title of the first chapter}

\section{Title of the first subchapter in the first chapter}

\section{Title of the second subchapter in the first chapter }


\chapter{Title of the second chapter}

\section{Title of the first subchapter in the second chapter}

\section{Title of the second subchapter in the second chapter}



% Ukázka použití některých konstrukcí LateXu (odkomentujte, chcete-li)
%%% Ukázka použití některých konstrukcí LaTeXu

\begin{otherlanguage}{czech}

\subsection{Ukázka \LaTeX{}u}
\label{ssec:ukazka}

V~této krátké části ukážeme použití několika základních konstrukcí \LaTeX{}u,
které by se vám mohly při psaní práce hodit.

Třeba odrážky:

\begin{itemize}
\item Logo Matfyzu vidíme na obrázku~\ref{fig:mff}.
\item Tato subsekce má číslo~\ref{ssec:ukazka}.
\item Odkaz na literaturu~\cite{lamport94}.
\end{itemize}

Druhy pomlček:
červeno-černý (krátká),
strana 16--22 (střední),
$45-44$ (minus),
a~toto je --- jak se asi dalo čekat --- vložená věta ohraničená dlouhými pomlčkami.
(Všimněte si, že jsme za \verb|a| napsali vlnovku místo mezery: to aby se
tam nemohl rozdělit řádek.)

% Makro na české uvozovky (novější verze LaTeXu ho už mají zabudované)
%\newcommand{\uv}[1]{\quotedblbase #1\textquotedblleft}
\uv{České uvozovky.}

\newtheorem{theorem}{Věta}
\newtheorem*{define}{Definice}	% Definice nečíslujeme, proto "*"

\begin{define}
{\sl Strom} je souvislý graf bez kružnic.
\end{define}

\begin{theorem}
Tato věta neplatí.
\end{theorem}

\begin{proof}
Neplatné věty nemají důkaz.
\end{proof}

\begin{figure}
	\centering
	\includegraphics[width=30mm]{img/logo.eps}
	\caption{Logo MFF UK}
	\label{fig:mff}
\end{figure}

\end{otherlanguage}


\chapter*{Conclusion}
\addcontentsline{toc}{chapter}{Conclusion}


%%% Seznam použité literatury
%%% Seznam použité literatury je zpracován podle platných standardů. Povinnou citační
%%% normou pro bakalářskou práci je ISO 690. Jména časopisů lze uvádět zkráceně, ale jen
%%% v kodifikované podobě. Všechny použité zdroje a prameny musí být řádně citovány.

\def\bibname{Bibliography}
\begin{thebibliography}{99}
\addcontentsline{toc}{chapter}{\bibname}

%\bibitem{lamport94}
%  {\sc Lamport,} Leslie.
%  \emph{\LaTeX: A Document Preparation System}.
%  2. vydání.
%  Massachusetts: Addison Wesley, 1994.
%  ISBN 0-201-52983-1.

\end{thebibliography}


%%% Tabulky v bakalářské práci, existují-li.
\chapwithtoc{List of Tables}

%%% Použité zkratky v bakalářské práci, existují-li, včetně jejich vysvětlení.
\chapwithtoc{List of Abbreviations}

%%% Přílohy k bakalářské práci, existují-li (různé dodatky jako výpisy
%%% programů, diagramy apod.). Každá příloha musí být alespoň jednou
%%% odkazována z vlastního textu práce. Přílohy se číslují.
\chapwithtoc{Attachments}

\openright
\end{document}
